\documentclass[]{article}

%opening
\title{The role of facilitation on evolutionary rescue in semi-arid environments}
\author{}

\begin{document}

\maketitle

\begin{abstract}
	
The capacity of a population to respond to a changing climate is determined by a complex balance between local adaptation and dispersal.
Most theoretical research on evolutionary rescue has focused on populations adapted to different environments in different places.
Here, we argue that ecological facilitation, i.e. the process by which some species facilitate the existence of others in harsh environments, is a special case of local adaptation to a heterogeneous environment with good and bad patches for the facilitated species, and that in arid environments this mimics habitat deterioration undergone by species under a climate change scenario. 
We provide a toy model to show how results from theory on local adaptation in heterogeneous environments and evolutionary rescue can help understand the adaptive dynamics of facilitated plants in semi-arid environments.

\end{abstract}

\section*{Methods}

Assume a landscape consisting of $N$ patches, all connected to each other by dispersal (this is an island model where all patches are equally distant).
The model unfolds in discrete time and generations are non-overlapping.
Each patch hosts a population of plants undergoing the following life cycle.

\end{document}
